% --- Archon physics formalization source begin ---
% archon:physics
% archon:covers PhyXMiniProblems/problem_phyx_mini_0259.lean
% archon:source-report reports/phyx_mini/problem_phyx_mini_0259.source.json

\chapter{Physics problem phyx\_mini\_0259}
\label{ch:PhyXMiniProblems_problem_phyx_mini_0259}

\paragraph{Problem source.}
\#\# Physical scenario

Figure shows the kinetic energy \$K\$ of a simple harmonic oscillator versus its position \$x\$. The vertical axis scale is set by \$K\_s = 4.0\textbackslash{} J\$.

\#\# Question

What is the spring constant?

\#\# Answer choices

- A: A: \$7.9 \textbackslash{}times 10\textasciicircum{}2\textbackslash{}\textbackslash{} N/m\$

- B: B: \$8.0 \textbackslash{}times 10\textasciicircum{}2\textbackslash{}\textbackslash{} N/m\$

- C: C: \$8.3 \textbackslash{}times 10\textasciicircum{}2\textbackslash{}\textbackslash{} N/m\$

- D: D: \$8.6 \textbackslash{}times 10\textasciicircum{}2\textbackslash{}\textbackslash{} N/m\$

\#\# Figure caption (auxiliary; use the image as primary evidence)

The image is a graph depicting a parabolic curve. The horizontal axis is labeled as \textbackslash{}( x \textbackslash{}) (cm), with a range from \textbackslash{}(-12\textbackslash{}) to \textbackslash{}(12\textbackslash{}), marked in intervals of \textbackslash{}(4\textbackslash{}). The vertical axis is labeled as \textbackslash{}( K \textbackslash{}) (J), representing energy in joules. The graph shows a thick, inverted parabola, indicating a symmetric peak at \textbackslash{}( x = 0 \textbackslash{}).

The peak of the parabola does not specify a numerical value but is marked with the text \textbackslash{}( K\_s \textbackslash{}). This label, \textbackslash{}( K\_s \textbackslash{}), is indicated along the vertical axis, representing some specific energy value at the vertex of the parabola. The parabola extends symmetrically from the peak, decreasing as it moves away from the center on both sides toward the x-axis at the endpoints. The grid is visible, providing a sense of scale, but the specific numerical values on the vertical axis are not shown apart from \textbackslash{}( K\_s \textbackslash{}).

\#\# Dataset metadata

Resonance and Harmonics; Physical Model Grounding Reasoning, Multi-Formula Reasoning

\paragraph{Recorded answer/context.}
C: C: \$8.3 \textbackslash{}times 10\textasciicircum{}2\textbackslash{}\textbackslash{} N/m\$

\paragraph{Figure/image path.}
/root/proposal\_for\_physic/science-mango/phyx\_mini\_run/phyx\_data/test\_image/259.png

\paragraph{Formalization target.}
create a compiling Lean file with sorry bodies at `PhyXMiniProblems/problem\_phyx\_mini\_0259.lean`. The Lean declarations must preserve the physical quantities, dimensions or dimensional roles, figure labels, governing-law hypotheses, and final relation expressed by this problem.
Use Mathlib/Physlib names found through LeanExplore where available. If a physics API is missing, introduce faithful local abstractions rather than scalar placeholder aliases.

\begin{theorem}[Physics formalization target]
\label{thm:physics:phyx_mini_0259:target}
\lean{PhyXMiniProblems.ProblemPhyXMini0259.problem_phyx_mini_0259}
\uses{def:physics:phyx-mini-0259:phyxminiproblems-problemphyxmini0259-stiffnessinnewtonspermeter, def:physics:phyx-mini-0259:phyxminiproblems-problemphyxmini0259-kineticenergygraphsetup, def:physics:phyx-mini-0259:phyxminiproblems-problemphyxmini0259-matchesprimaryfigure, def:physics:phyx-mini-0259:phyxminiproblems-problemphyxmini0259-matchesproblemdata, def:physics:phyx-mini-0259:phyxminiproblems-problemphyxmini0259-hasphysicalparameters, def:physics:phyx-mini-0259:phyxminiproblems-problemphyxmini0259-satisfiessimpleharmonicenergylaw, def:physics:phyx-mini-0259:phyxminiproblems-problemphyxmini0259-recordedanswerchoice, def:physics:phyx-mini-0259:phyxminiproblems-problemphyxmini0259-matchesanswerchoice, def:physics:phyx-mini-0259:phyxminiproblems-problemphyxmini0259-isuniqueclosestanswerchoice}
The assigned autoformalize agent should translate this physics problem into Lean declarations in the covered file, with theorem and lemma proof bodies written as `by sorry`.
\end{theorem}
\begin{proof}
This is an autoformalization task, not a proof task. Produce statements that can later be proved by the physics prover without weakening the physical model.
\end{proof}

% --- Archon declaration topology begin ---
\section{Declaration topology}
\begin{definition}[Length Quantity]
  \label{def:physics:phyx-mini-0259:phyxminiproblems-problemphyxmini0259-lengthquantity}
  \lean{PhyXMiniProblems.ProblemPhyXMini0259.LengthQuantity}
  A signed displacement from the oscillator's equilibrium position
\end{definition}
\begin{proof}
  This is the declared model definition of Length Quantity.
\end{proof}
\begin{definition}[Mass Quantity]
  \label{def:physics:phyx-mini-0259:phyxminiproblems-problemphyxmini0259-massquantity}
  \lean{PhyXMiniProblems.ProblemPhyXMini0259.MassQuantity}
  The physical mass of the oscillating body
\end{definition}
\begin{proof}
  This is the declared model definition of Mass Quantity.
\end{proof}
\begin{definition}[Spring Stiffness Quantity]
  \label{def:physics:phyx-mini-0259:phyxminiproblems-problemphyxmini0259-springstiffnessquantity}
  \lean{PhyXMiniProblems.ProblemPhyXMini0259.SpringStiffnessQuantity}
  Spring stiffness, whose dimension is force per length, M T⁻²
\end{definition}
\begin{proof}
  This is the declared model definition of Spring Stiffness Quantity.
\end{proof}
\begin{definition}[Energy Quantity]
  \label{def:physics:phyx-mini-0259:phyxminiproblems-problemphyxmini0259-energyquantity}
  \lean{PhyXMiniProblems.ProblemPhyXMini0259.EnergyQuantity}
  Mechanical energy, whose dimension is M L² T⁻²
\end{definition}
\begin{proof}
  This is the declared model definition of Energy Quantity.
\end{proof}
\begin{definition}[centimeter Unit Choices]
  \label{def:physics:phyx-mini-0259:phyxminiproblems-problemphyxmini0259-centimeterunitchoices}
  \lean{PhyXMiniProblems.ProblemPhyXMini0259.centimeterUnitChoices}
  SI base units with centimetres selected as the length unit
\end{definition}
\begin{proof}
  This is the declared model definition of centimeter Unit Choices.
\end{proof}
\begin{definition}[length In Meters]
  \label{def:physics:phyx-mini-0259:phyxminiproblems-problemphyxmini0259-lengthinmeters}
  \lean{PhyXMiniProblems.ProblemPhyXMini0259.lengthInMeters}
  \uses{def:physics:phyx-mini-0259:phyxminiproblems-problemphyxmini0259-lengthquantity}
  Scalar readout of a signed physical displacement in metres
\end{definition}
\begin{proof}
  This is the declared model definition of length In Meters.
\end{proof}
\begin{definition}[length In Centimeters]
  \label{def:physics:phyx-mini-0259:phyxminiproblems-problemphyxmini0259-lengthincentimeters}
  \lean{PhyXMiniProblems.ProblemPhyXMini0259.lengthInCentimeters}
  \uses{def:physics:phyx-mini-0259:phyxminiproblems-problemphyxmini0259-lengthquantity, def:physics:phyx-mini-0259:phyxminiproblems-problemphyxmini0259-centimeterunitchoices}
  Scalar readout of a signed physical displacement in centimetres
\end{definition}
\begin{proof}
  This is the declared model definition of length In Centimeters.
\end{proof}
\begin{definition}[mass In Kilograms]
  \label{def:physics:phyx-mini-0259:phyxminiproblems-problemphyxmini0259-massinkilograms}
  \lean{PhyXMiniProblems.ProblemPhyXMini0259.massInKilograms}
  \uses{def:physics:phyx-mini-0259:phyxminiproblems-problemphyxmini0259-massquantity}
  Scalar readout of a physical mass in kilograms
\end{definition}
\begin{proof}
  This is the declared model definition of mass In Kilograms.
\end{proof}
\begin{definition}[stiffness In Newtons Per Meter]
  \label{def:physics:phyx-mini-0259:phyxminiproblems-problemphyxmini0259-stiffnessinnewtonspermeter}
  \lean{PhyXMiniProblems.ProblemPhyXMini0259.stiffnessInNewtonsPerMeter}
  \uses{def:physics:phyx-mini-0259:phyxminiproblems-problemphyxmini0259-springstiffnessquantity}
  Scalar readout of spring stiffness in newtons per metre
\end{definition}
\begin{proof}
  This is the declared model definition of stiffness In Newtons Per Meter.
\end{proof}
\begin{definition}[energy In Joules]
  \label{def:physics:phyx-mini-0259:phyxminiproblems-problemphyxmini0259-energyinjoules}
  \lean{PhyXMiniProblems.ProblemPhyXMini0259.energyInJoules}
  \uses{def:physics:phyx-mini-0259:phyxminiproblems-problemphyxmini0259-energyquantity}
  Scalar readout of a mechanical energy in joules
\end{definition}
\begin{proof}
  This is the declared model definition of energy In Joules.
\end{proof}
\begin{definition}[Graph Axis]
  \label{def:physics:phyx-mini-0259:phyxminiproblems-problemphyxmini0259-graphaxis}
  \lean{PhyXMiniProblems.ProblemPhyXMini0259.GraphAxis}
  The two coordinate axes visible in the supplied graph
\end{definition}
\begin{proof}
  This is the declared model definition of Graph Axis.
\end{proof}
\begin{definition}[Axis Observable]
  \label{def:physics:phyx-mini-0259:phyxminiproblems-problemphyxmini0259-axisobservable}
  \lean{PhyXMiniProblems.ProblemPhyXMini0259.AxisObservable}
  The physical observable plotted on each graph axis
\end{definition}
\begin{proof}
  This is the declared model definition of Axis Observable.
\end{proof}
\begin{definition}[Figure Label]
  \label{def:physics:phyx-mini-0259:phyxminiproblems-problemphyxmini0259-figurelabel}
  \lean{PhyXMiniProblems.ProblemPhyXMini0259.FigureLabel}
  Text labels printed in the graph
\end{definition}
\begin{proof}
  This is the declared model definition of Figure Label.
\end{proof}
\begin{definition}[Figure Feature]
  \label{def:physics:phyx-mini-0259:phyxminiproblems-problemphyxmini0259-figurefeature}
  \lean{PhyXMiniProblems.ProblemPhyXMini0259.FigureFeature}
  Graph features to which the printed labels refer
\end{definition}
\begin{proof}
  This is the declared model definition of Figure Feature.
\end{proof}
\begin{definition}[Horizontal Tick]
  \label{def:physics:phyx-mini-0259:phyxminiproblems-problemphyxmini0259-horizontaltick}
  \lean{PhyXMiniProblems.ProblemPhyXMini0259.HorizontalTick}
  The seven numbered horizontal-axis ticks in the primary image
\end{definition}
\begin{proof}
  This is the declared model definition of Horizontal Tick.
\end{proof}
\begin{definition}[centimeter Readout]
  \label{def:physics:phyx-mini-0259:phyxminiproblems-problemphyxmini0259-horizontaltick-centimeterreadout}
  \lean{PhyXMiniProblems.ProblemPhyXMini0259.HorizontalTick.centimeterReadout}
  \uses{def:physics:phyx-mini-0259:phyxminiproblems-problemphyxmini0259-horizontaltick}
  Centimetre number printed at a horizontal-axis tick
\end{definition}
\begin{proof}
  This is the declared model definition of centimeter Readout.
\end{proof}
\begin{definition}[Kinetic Energy Graph Setup]
  \label{def:physics:phyx-mini-0259:phyxminiproblems-problemphyxmini0259-kineticenergygraphsetup}
  \lean{PhyXMiniProblems.ProblemPhyXMini0259.KineticEnergyGraphSetup}
  \uses{def:physics:phyx-mini-0259:phyxminiproblems-problemphyxmini0259-lengthquantity, def:physics:phyx-mini-0259:phyxminiproblems-problemphyxmini0259-massquantity, def:physics:phyx-mini-0259:phyxminiproblems-problemphyxmini0259-springstiffnessquantity, def:physics:phyx-mini-0259:phyxminiproblems-problemphyxmini0259-energyquantity, def:physics:phyx-mini-0259:phyxminiproblems-problemphyxmini0259-graphaxis, def:physics:phyx-mini-0259:phyxminiproblems-problemphyxmini0259-axisobservable, def:physics:phyx-mini-0259:phyxminiproblems-problemphyxmini0259-figurelabel, def:physics:phyx-mini-0259:phyxminiproblems-problemphyxmini0259-figurefeature, def:physics:phyx-mini-0259:phyxminiproblems-problemphyxmini0259-horizontaltick}
  The physical oscillator and the independently stored graph observables. kineticEnergyAtDisplacement is the measured curve rather than a definition from the requested stiffness. oscillatorSIReadout and positionVectorInMeters are the scalar data needed to use Physlib's one-dimensional harmonic-oscillator potential energy. No numerical value for the requested stiffness is assigned here
\end{definition}
\begin{proof}
  This is the declared model definition of Kinetic Energy Graph Setup.
\end{proof}
\begin{definition}[Is In Displayed Domain]
  \label{def:physics:phyx-mini-0259:phyxminiproblems-problemphyxmini0259-isindisplayeddomain}
  \lean{PhyXMiniProblems.ProblemPhyXMini0259.IsInDisplayedDomain}
  \uses{def:physics:phyx-mini-0259:phyxminiproblems-problemphyxmini0259-lengthquantity, def:physics:phyx-mini-0259:phyxminiproblems-problemphyxmini0259-lengthinmeters, def:physics:phyx-mini-0259:phyxminiproblems-problemphyxmini0259-kineticenergygraphsetup}
  A displacement lies in the physical interval displayed between the two turning points of the graph
\end{definition}
\begin{proof}
  This is the declared model definition of Is In Displayed Domain.
\end{proof}
\begin{definition}[Matches Primary Figure]
  \label{def:physics:phyx-mini-0259:phyxminiproblems-problemphyxmini0259-matchesprimaryfigure}
  \lean{PhyXMiniProblems.ProblemPhyXMini0259.MatchesPrimaryFigure}
  \uses{def:physics:phyx-mini-0259:phyxminiproblems-problemphyxmini0259-lengthincentimeters, def:physics:phyx-mini-0259:phyxminiproblems-problemphyxmini0259-energyinjoules, def:physics:phyx-mini-0259:phyxminiproblems-problemphyxmini0259-kineticenergygraphsetup, def:physics:phyx-mini-0259:phyxminiproblems-problemphyxmini0259-isindisplayeddomain}
  Qualitative geometry and quantitative readouts taken from the primary bitmap. The K\_s mark is four vertical grid steps above zero, while the peak at x = 0 is six steps above zero. The ratio field records this graphical scale comparison without assuming any value for the spring stiffness
\end{definition}
\begin{proof}
  This is the declared model definition of Matches Primary Figure.
\end{proof}
\begin{definition}[Matches Problem Data]
  \label{def:physics:phyx-mini-0259:phyxminiproblems-problemphyxmini0259-matchesproblemdata}
  \lean{PhyXMiniProblems.ProblemPhyXMini0259.MatchesProblemData}
  \uses{def:physics:phyx-mini-0259:phyxminiproblems-problemphyxmini0259-energyinjoules, def:physics:phyx-mini-0259:phyxminiproblems-problemphyxmini0259-kineticenergygraphsetup}
  The numerical vertical scale stated with the problem. The peak kinetic energy and the spring stiffness remain unconstrained here
\end{definition}
\begin{proof}
  This is the declared model definition of Matches Problem Data.
\end{proof}
\begin{definition}[Has Physical Parameters]
  \label{def:physics:phyx-mini-0259:phyxminiproblems-problemphyxmini0259-hasphysicalparameters}
  \lean{PhyXMiniProblems.ProblemPhyXMini0259.HasPhysicalParameters}
  \uses{def:physics:phyx-mini-0259:phyxminiproblems-problemphyxmini0259-lengthinmeters, def:physics:phyx-mini-0259:phyxminiproblems-problemphyxmini0259-massinkilograms, def:physics:phyx-mini-0259:phyxminiproblems-problemphyxmini0259-stiffnessinnewtonspermeter, def:physics:phyx-mini-0259:phyxminiproblems-problemphyxmini0259-energyinjoules, def:physics:phyx-mini-0259:phyxminiproblems-problemphyxmini0259-kineticenergygraphsetup, def:physics:phyx-mini-0259:phyxminiproblems-problemphyxmini0259-isindisplayeddomain}
  Positivity and nondegeneracy conditions for the physical oscillator and the displayed kinetic-energy curve
\end{definition}
\begin{proof}
  This is the declared model definition of Has Physical Parameters.
\end{proof}
\begin{definition}[Satisfies Simple Harmonic Energy Law]
  \label{def:physics:phyx-mini-0259:phyxminiproblems-problemphyxmini0259-satisfiessimpleharmonicenergylaw}
  \lean{PhyXMiniProblems.ProblemPhyXMini0259.SatisfiesSimpleHarmonicEnergyLaw}
  \uses{def:physics:phyx-mini-0259:phyxminiproblems-problemphyxmini0259-lengthinmeters, def:physics:phyx-mini-0259:phyxminiproblems-problemphyxmini0259-massinkilograms, def:physics:phyx-mini-0259:phyxminiproblems-problemphyxmini0259-stiffnessinnewtonspermeter, def:physics:phyx-mini-0259:phyxminiproblems-problemphyxmini0259-energyinjoules, def:physics:phyx-mini-0259:phyxminiproblems-problemphyxmini0259-kineticenergygraphsetup, def:physics:phyx-mini-0259:phyxminiproblems-problemphyxmini0259-isindisplayeddomain}
  The governing simple-harmonic-oscillator energy model in coherent SI readouts. Physlib's scalar oscillator is linked to the dimensionful physical mass and stiffness. Its potentialEnergy supplies U(x) = (1/2) k x², and conservation of mechanical energy gives K(x) + U(x) = K(0) + U(0) throughout the accessible interval. These are general model laws: no target stiffness or answer choice appears in this interface
\end{definition}
\begin{proof}
  This is the declared model definition of Satisfies Simple Harmonic Energy Law.
\end{proof}
\begin{definition}[Answer Choice]
  \label{def:physics:phyx-mini-0259:phyxminiproblems-problemphyxmini0259-answerchoice}
  \lean{PhyXMiniProblems.ProblemPhyXMini0259.AnswerChoice}
  Labels printed beside the four candidate spring constants
\end{definition}
\begin{proof}
  This is the declared model definition of Answer Choice.
\end{proof}
\begin{definition}[newtons Per Meter]
  \label{def:physics:phyx-mini-0259:phyxminiproblems-problemphyxmini0259-answerchoice-newtonspermeter}
  \lean{PhyXMiniProblems.ProblemPhyXMini0259.AnswerChoice.newtonsPerMeter}
  \uses{def:physics:phyx-mini-0259:phyxminiproblems-problemphyxmini0259-answerchoice}
  Spring stiffness printed beside an answer, in newtons per metre
\end{definition}
\begin{proof}
  This is the declared model definition of newtons Per Meter.
\end{proof}
\begin{definition}[recorded Answer Choice]
  \label{def:physics:phyx-mini-0259:phyxminiproblems-problemphyxmini0259-recordedanswerchoice}
  \lean{PhyXMiniProblems.ProblemPhyXMini0259.recordedAnswerChoice}
  \uses{def:physics:phyx-mini-0259:phyxminiproblems-problemphyxmini0259-answerchoice}
  The answer label recorded in the supplied dataset metadata
\end{definition}
\begin{proof}
  This is the declared model definition of recorded Answer Choice.
\end{proof}
\begin{definition}[Matches Answer Choice]
  \label{def:physics:phyx-mini-0259:phyxminiproblems-problemphyxmini0259-matchesanswerchoice}
  \lean{PhyXMiniProblems.ProblemPhyXMini0259.MatchesAnswerChoice}
  \uses{def:physics:phyx-mini-0259:phyxminiproblems-problemphyxmini0259-stiffnessinnewtonspermeter, def:physics:phyx-mini-0259:phyxminiproblems-problemphyxmini0259-kineticenergygraphsetup, def:physics:phyx-mini-0259:phyxminiproblems-problemphyxmini0259-answerchoice}
  Agreement with a value displayed to the nearest ten newtons per metre
\end{definition}
\begin{proof}
  This is the declared model definition of Matches Answer Choice.
\end{proof}
\begin{definition}[Is Unique Closest Answer Choice]
  \label{def:physics:phyx-mini-0259:phyxminiproblems-problemphyxmini0259-isuniqueclosestanswerchoice}
  \lean{PhyXMiniProblems.ProblemPhyXMini0259.IsUniqueClosestAnswerChoice}
  \uses{def:physics:phyx-mini-0259:phyxminiproblems-problemphyxmini0259-stiffnessinnewtonspermeter, def:physics:phyx-mini-0259:phyxminiproblems-problemphyxmini0259-kineticenergygraphsetup, def:physics:phyx-mini-0259:phyxminiproblems-problemphyxmini0259-answerchoice}
  The selected displayed value is strictly closer than every other option
\end{definition}
\begin{proof}
  This is the declared model definition of Is Unique Closest Answer Choice.
\end{proof}
\begin{lemma}[spring Stiffness exact]
  \label{lem:physics:phyx-mini-0259:phyxminiproblems-problemphyxmini0259-springstiffness-exact}
  \lean{PhyXMiniProblems.ProblemPhyXMini0259.springStiffness_exact}
  \uses{def:physics:phyx-mini-0259:phyxminiproblems-problemphyxmini0259-stiffnessinnewtonspermeter, def:physics:phyx-mini-0259:phyxminiproblems-problemphyxmini0259-kineticenergygraphsetup, def:physics:phyx-mini-0259:phyxminiproblems-problemphyxmini0259-matchesprimaryfigure, def:physics:phyx-mini-0259:phyxminiproblems-problemphyxmini0259-matchesproblemdata, def:physics:phyx-mini-0259:phyxminiproblems-problemphyxmini0259-hasphysicalparameters, def:physics:phyx-mini-0259:phyxminiproblems-problemphyxmini0259-satisfiessimpleharmonicenergylaw}
  The graph peak is (6 / 4) K\_s = 6 J and the turning-point magnitude is 0.12 m. Conservation with U(A) = (1/2) k A² therefore gives the exact SI readout k = 2500/3 N/m
\end{lemma}
\begin{proof}
  The conclusion follows from the governing relations and figure readouts named in the statement of spring Stiffness exact.
\end{proof}
% --- Archon declaration topology end ---
% --- Archon physics formalization source end ---
